\Article{39}
Um deutlicher zu erkennen, mit welchem Recht der zufällige Wert von \(M\) seinem Mittelwert gleichgesetzt werden darf, ist noch der mittlere Fehler zu untersuchen, der zu befürchten ist, wenn wir \(\frac{M}{\pi-\rho}=\mu\mu\) setzen. Dieser mittlere Fehler ist gleich der Quadratwurzel aus dem Mittelwert der Größe
\[
\left(\frac{\Omega^0-x^0\xi^0-y^0\eta^0-z^0\zeta^0-\text{etc.}-(\pi-\rho)\mu\mu}{\pi-\rho}\right)^2
\]
die wir so darstellen werden:
\[
\begin{aligned}
&\left(\frac{\Omega^0-x^0\xi^0-y^0\eta^0-z^0\zeta^0-\text{etc.}}{\pi-\rho}\right)^2\\[4pt]
&\quad-\frac{2\mu\mu}{\pi-\rho}(\Omega^0-x^0\xi^0-y^0\eta^0-z^0\zeta^0-\text{etc.}-(\pi-\rho)\mu\mu)-\mu^4
\end{aligned}
\]
Da offenbar der Mittelwert des zweiten Glieds \(=0\) ist, kommt es darauf an, dass wir den Mittelwert der Funktion
\[
\Psi=(\Omega^0-x^0\xi^0-y^0\eta^0-z^0\zeta^0-\text{etc.})^2
\]
ermitteln. Ist dieser gefunden und mit \(N\) bezeichnet, so ist der gesuchte mittlere Fehler
\[
=\BareRoot{\left(\frac{N}{(\pi-\rho)^2}-\mu^4\right)}
\]
Der entwickelte Ausdruck \(\Psi\) ist offenbar eine homogene Funktion entweder der Fehler \(e,e',e''\) etc. oder der Größen \(\varepsilon,\varepsilon',\varepsilon''\) etc.; sein Mittelwert wird gefunden, wenn

\(1^\mathrm{o}\) anstelle der vierten Potenzen \(e^4,e'^4,e''^4\) etc. deren Mittelwerte eingesetzt werden,

\(2^\mathrm{o}\) anstelle jedes Produkts aus zwei Quadraten, etwa \(eee'e',\allowbreak eee''e'',\allowbreak e'e'e''e''\) etc., die Produkte ihrer Mittelwerte, nämlich \(mmm'm',\allowbreak mmm''m'',\allowbreak m'm'm''m''\) etc., eingesetzt werden,
\par\SourceLeaf{055}%
\(3^\mathrm{o}\) die übrigen Teile aber, die entweder einen Faktor wie \(e^3e'\) oder wie \(eee'e''\) enthalten werden, gänzlich weggelassen werden.

Wir nehmen an, dass die Mittelwerte der vierten Potenzen \(e^4,\allowbreak e'^4,\allowbreak e''^4\) etc. den vierten Potenzen \(m^4,\allowbreak m'^4,\allowbreak m''^4\) etc. selbst proportional sind (siehe Art.\ 16), sodass jene sich zu diesen wie \(\nu^4\) zu \(\mu^4\) verhalten und somit \(\nu^4\) den Mittelwert der vierten Potenzen solcher Beobachtungen bezeichnet, deren Gewicht \(=1\) ist. Daher können die vorhergehenden Vorschriften auch so ausgedrückt werden: Anstelle jeder vierten Potenz \(\varepsilon^4,\varepsilon'^4,\varepsilon''^4\) etc. ist \(\nu^4\) zu schreiben, anstelle jedes Produkts aus zwei Quadraten, etwa \(\varepsilon\varepsilon\varepsilon'\varepsilon',\allowbreak\varepsilon\varepsilon\varepsilon''\varepsilon'',\allowbreak\varepsilon'\varepsilon'\varepsilon''\varepsilon''\) etc., ist \(\mu^4\) zu schreiben; alle übrigen Glieder, die Faktoren wie \(\varepsilon^3\varepsilon'\) oder \(\varepsilon\varepsilon\varepsilon'\varepsilon''\) oder \(\varepsilon\varepsilon'\varepsilon''\varepsilon'''\) enthalten werden, sind wegzulassen.

Ist dies richtig verstanden, so wird leicht ersichtlich:

I. Der Mittelwert des Quadrats \(\Omega^0\Omega^0\) ist \(\pi\nu^4+(\pi\pi-\pi)\mu^4\)

II. Der Mittelwert des Produkts \(\varepsilon\varepsilon x^0\xi^0\) ist \(=a\alpha\nu^4+(a'\alpha'+a''\alpha''+\text{etc.})\mu^4\) oder wegen \(a\alpha+a'\alpha'+a''\alpha''+\text{etc.}=1\)
\[
=a\alpha(\nu^4-\mu^4)+\mu^4
\]
Da ebenso der Mittelwert des Produkts \(\varepsilon'\varepsilon'x^0\xi^0\) gleich \(=a'\alpha'(\nu^4-\mu^4)+\mu^4\) ist, der Mittelwert des Produkts \(\varepsilon''\varepsilon''x^0\xi^0\) aber gleich \(=a''\alpha''(\nu^4-\mu^4)+\mu^4\) und so fort, ist klar, dass der Mittelwert des Produkts \((\varepsilon\varepsilon+\varepsilon'\varepsilon'+\varepsilon''\varepsilon''+\text{etc.})x^0\xi^0\) oder \(\Omega^0x^0\xi^0\) gleich dem folgenden Ausdruck ist:
\[
=\nu^4-\mu^4+\pi\mu^4
\]
\begingroup\setlength{\rightskip}{0pt plus 2em}\noindent Denselben Mittelwert haben die Produkte \(\Omega^0y^0\eta^0,\allowbreak\Omega^0z^0\zeta^0\) etc. Deshalb ist der Mittelwert des Produkts \(\Omega^0(x^0\xi^0+y^0\eta^0+z^0\zeta^0+\text{etc.})\) gleich dem folgenden Ausdruck:\par\endgroup
\[
=\rho\nu^4+\rho(\pi-1)\mu^4
\]

III. Damit die übrigen Entwicklungen nicht zu weitläufig werden, ist eine geeignete Bezeichnung einzuführen. Wir verwenden daher das Zeichen \(\Sigma\) in einem etwas weiteren Sinn, als es oben an verschiedenen Stellen geschehen ist: Es bezeichne die Summe des Glieds, dem es vorangestellt wird, und aller ähnlichen, aber nicht identischen Glieder, die daraus durch sämtliche Vertauschungen der Beobachtungen entstehen. So haben wir z.\ B. \(x^0=\Sigma\alpha\varepsilon\), \(x^0x^0=\Sigma\alpha\alpha\varepsilon\varepsilon+2\Sigma\alpha\alpha'\varepsilon\varepsilon'\). Stellen wir den Mittelwert des Produkts \(x^0x^0\xi^0\xi^0\) nach seinen Teilen zusammen, so haben wir zunächst den Mittelwert des Produkts \(\alpha\alpha\varepsilon\varepsilon\xi^0\xi^0\)
\[
\begin{aligned}
&=aa\alpha\alpha\nu^4+\alpha\alpha(a'a'+a''a''+\text{etc.})\mu^4\\
&=aa\alpha\alpha(\nu^4-\mu^4)+\alpha\alpha\mu^4\Sigma aa
\end{aligned}
\]
\SourceLeaf{056}
Ebenso ist der Mittelwert des Produkts \(\alpha'\alpha'\varepsilon'\varepsilon'\xi^0\xi^0\) gleich \(=a'a'\alpha'\alpha'(\nu^4-\mu^4)+\alpha'\alpha'\mu^4\Sigma aa\) und so fort; daher ist der Mittelwert des Produkts \(\xi^0\xi^0\Sigma\alpha\alpha\varepsilon\varepsilon\) gleich dem folgenden Ausdruck:
\[
=(\nu^4-\mu^4)\Sigma aa\alpha\alpha+\mu^4\Sigma aa.\Sigma\alpha\alpha
\]
Ferner ist der Mittelwert des Produkts \(\alpha\alpha'\varepsilon\varepsilon'\xi^0\xi^0\) gleich \(=2aa'\alpha\alpha'\mu^4\), der Mittelwert des Produkts \(\alpha\alpha''\varepsilon\varepsilon''\xi^0\xi^0\) ebenso gleich \(=2aa''\alpha\alpha''\mu^4\) etc. Daraus wird leicht geschlossen, dass der Mittelwert des Produkts \(\xi^0\xi^0\Sigma\alpha\alpha'\varepsilon\varepsilon'\) gleich dem folgenden Ausdruck ist:
\[
\begin{aligned}
&=2\mu^4\Sigma a\alpha a'\alpha'=\mu^4((\Sigma a\alpha)^2-\Sigma aa\alpha\alpha)\\
&=\mu^4(1-\Sigma aa\alpha\alpha)
\end{aligned}
\]
Fassen wir dies zusammen, so haben wir den Mittelwert des Produkts \(x^0x^0\xi^0\xi^0\)
\[
=(\nu^4-3\mu^4)\Sigma aa\alpha\alpha+2\mu^4+\mu^4\Sigma aa.\Sigma\alpha\alpha
\]

IV. Auf ganz ähnliche Weise findet sich für den Mittelwert des Produkts \(x^0y^0\xi^0\eta^0\) der Ausdruck
\[
\begin{aligned}
&=\nu^4\Sigma ab\alpha\beta+\mu^4\Sigma a\alpha b'\beta'\\
&\quad+\mu^4\Sigma ab\alpha'\beta'+\mu^4\Sigma a\beta b'\alpha'
\end{aligned}
\]
Es gilt aber
\[
\begin{aligned}
\Sigma a\alpha b'\beta'&=\Sigma a\alpha.\Sigma b\beta-\Sigma a\alpha b\beta\\
\Sigma ab\alpha'\beta'&=\Sigma ab.\Sigma\alpha\beta-\Sigma ab\alpha\beta\\
\Sigma a\beta b'\alpha'&=\Sigma a\beta.\Sigma b\alpha-\Sigma a\beta b\alpha
\end{aligned}
\]
sodass jener Mittelwert wegen \(\Sigma a\alpha=1,\allowbreak\Sigma b\beta=1,\allowbreak\Sigma a\beta=0,\allowbreak\Sigma b\alpha=0\) den folgenden Ausdruck erhält:
\[
=(\nu^4-3\mu^4)\Sigma ab\alpha\beta+\mu^4(1+\Sigma ab.\Sigma\alpha\beta)
\]

V. Da auf ganz dieselbe Weise der Mittelwert des Produkts \(x^0z^0\xi^0\zeta^0\) gleich dem folgenden Ausdruck ist:
\[
=(\nu^4-3\mu^4)\Sigma ac\alpha\gamma+\mu^4(1+\Sigma ac.\Sigma\alpha\gamma)
\]
und so fort, liefert die Addition für den Mittelwert des Produkts \(x^0\xi^0(x^0\xi^0+y^0\eta^0+z^0\zeta^0+\text{etc.})\) den Ausdruck
\[
\begin{aligned}
&=(\nu^4-3\mu^4)\Sigma(a\alpha(a\alpha+b\beta+c\gamma+\text{etc.}))+(\rho+1)\mu^4\\
&\quad+\mu^4(\Sigma aa.\Sigma\alpha\alpha+\Sigma ab.\Sigma\alpha\beta+\Sigma ac.\Sigma\alpha\gamma+\text{etc.})\\
&=(\nu^4-3\mu^4)\Sigma(a\alpha(a\alpha+b\beta+c\gamma+\text{etc.}))+(\rho+2)\mu^4
\end{aligned}
\]

VI. Auf ganz dieselbe Weise findet sich für den Mittelwert des Produkts \(y^0\eta^0(x^0\xi^0+y^0\eta^0+z^0\zeta^0+\text{etc.})\) der Ausdruck
\SourceLeaf{057}
\[
=(\nu^4-3\mu^4)\Sigma(b\beta(a\alpha+b\beta+c\gamma+\text{etc.}))+(\rho+2)\mu^4
\]
danach für den Mittelwert des Produkts \(z^0\zeta^0(x^0\xi^0+y^0\eta^0+z^0\zeta^0+\text{etc.})\) der Ausdruck
\[
=(\nu^4-3\mu^4)\Sigma(c\gamma(a\alpha+b\beta+c\gamma+\text{etc.}))+(\rho+2)\mu^4
\]
und so fort. Daraus ergibt sich durch Addition für den Mittelwert des Quadrats \((x^0\xi^0+y^0\eta^0+z^0\zeta^0+\text{etc.})^2\) der Ausdruck
\[
=(\nu^4-3\mu^4)\Sigma((a\alpha+b\beta+c\gamma+\text{etc.})^2)+(\rho\rho+2\rho)\mu^4
\]

VII. Fassen wir schließlich alles richtig zusammen, so ergibt sich
\[
\begin{aligned}
N&=(\pi-2\rho)\nu^4+(\pi\pi-\pi-2\pi\rho+4\rho+\rho\rho)\mu^4\\
&\quad+(\nu^4-3\mu^4)\Sigma((a\alpha+b\beta+c\gamma+\text{etc.})^2)\\
&=(\pi-\rho)(\nu^4-\mu^4)+(\pi-\rho)^2\mu^4\\
&\quad-(\nu^4-3\mu^4)[\rho-\Sigma((a\alpha+b\beta+c\gamma+\text{etc.})^2)]
\end{aligned}
\]
Der mittlere Fehler, der bei der Bestimmung von \(\mu\mu\) durch die Formel
\[
\mu\mu=\frac{M}{\pi-\rho}
\]
zu befürchten ist, ist daher
\[
=\BareRoot{\left\{\frac{\nu^4-\mu^4}{\pi-\rho}-\frac{\nu^4-3\mu^4}{(\pi-\rho)^2}.[\rho-\Sigma((a\alpha+b\beta+c\gamma+\text{etc.})^2)]\right\}}
\]

\Article{40}
Die Größe \(\Sigma((a\alpha+b\beta+c\gamma+\text{etc.})^2)\), die in den gerade gefundenen Ausdruck eingeht, lässt sich im Allgemeinen zwar nicht auf eine einfachere Form zurückführen; dennoch können zwei Grenzen angegeben werden, zwischen denen ihr Wert notwendig liegen muss. \textit{Erstens} lässt sich aus den oben entwickelten Beziehungen leicht beweisen, dass
\[
\begin{aligned}
&(a\alpha+b\beta+c\gamma+\text{etc.})^2+(a\alpha'+b\beta'+c\gamma'+\text{etc.})^2\\
&\quad+(a\alpha''+b\beta''+c\gamma''+\text{etc.})^2+\text{etc.}\\
&=a\alpha+b\beta+c\gamma+\text{etc.}
\end{aligned}
\]
woraus wir schließen, dass \(a\alpha+b\beta+c\gamma+\text{etc.}\) eine positive Größe ist, die kleiner als die Einheit ist (zumindest nicht größer). Dasselbe gilt für die Größe \(a'\alpha'+b'\beta'+c'\gamma'+\text{etc.}\), die der Summe
\[
\begin{aligned}
&(a'\alpha+b'\beta+c'\gamma+\text{etc.})^2+(a'\alpha'+b'\beta'+c'\gamma'+\text{etc.})^2\\
&\quad+(a'\alpha''+b'\beta''+c'\gamma''\text{ etc.})^2+\text{etc.}
\end{aligned}
\]
gleich gefunden wird; ebenso ist \(a''\alpha''+b''\beta''+c''\gamma''+\text{etc.}\) kleiner als die Einheit und so
\SourceLeaf{058}%
fort. Daher ist \(\Sigma((a\alpha+b\beta+c\gamma+\text{etc.})^2)\) notwendig kleiner als \(\pi\). \textit{Zweitens} gilt \(\Sigma(a\alpha+b\beta+c\gamma+\text{etc.})=\rho\), weil \(\Sigma a\alpha=1,\allowbreak\Sigma b\beta=1,\allowbreak\Sigma c\gamma=1\) etc. ist; daraus wird leicht abgeleitet, dass die Quadratsumme \(\Sigma((a\alpha+b\beta+c\gamma+\text{etc.})^2)\) größer als \(\frac{\rho\rho}{\pi}\) oder zumindest nicht kleiner ist. Daher liegt das Glied
\[
\frac{\nu^4-3\mu^4}{(\pi-\rho)^2}.[\rho-\Sigma((a\alpha+b\beta+c\gamma+\text{etc.})^2)]
\]
notwendig zwischen den Grenzen \(-\frac{\nu^4-3\mu^4}{\pi-\rho}\) und \(\frac{\nu^4-3\mu^4}{\pi-\rho}.\frac{\rho}{\pi}\) oder, wenn wir weitere vorziehen, zwischen \(-\frac{\nu^4-3\mu^4}{\pi-\rho}\) und \(+\frac{\nu^4-3\mu^4}{\pi-\rho}\); folglich liegt das Quadrat des mittleren Fehlers, der beim Wert von \(\mu\mu=\frac{M}{\pi-\rho}\) zu befürchten ist, zwischen den Grenzen \(\frac{2\nu^4-4\mu^4}{\pi-\rho}\) und \(\frac{2\mu^4}{\pi-\rho}\), sodass jede beliebige Genauigkeit erreicht werden kann, wenn nur die Zahl der Beobachtungen hinreichend groß ist.

Es ist sehr bemerkenswert, dass unter jener Hypothese (Art.\ 9, III), auf die die Theorie der kleinsten Quadrate früher gegründet war, jenes Glied vollständig wegfällt. Wie in allen Fällen zur Ermittlung eines Näherungswerts für den mittleren Beobachtungsfehler \(\mu\) die Summe \(\lambda\lambda+\lambda'\lambda'+\lambda''\lambda''+\text{etc.}=M\) so zu behandeln ist, als wäre sie die Summe von \(\pi-\rho\) zufälligen Fehlern, so wird unter jener Hypothese auch die Genauigkeit dieser Bestimmung selbst gleich derjenigen, die einer Bestimmung aus \(\pi-\rho\) wahren Fehlern zuzuschreiben ist, wie wir in Art.\ 15 gefunden haben.
\par\nobreak\medskip\centerline{\rule{.25\linewidth}{.3pt}}
\clearpage
